% !TEX program = xelatex
\documentclass[11pt,a4paper]{article}

\usepackage[margin=22mm,headheight=15pt]{geometry}
\usepackage{fontspec}
\usepackage{xeCJK}
\usepackage{amsmath,amssymb,mathtools}
\usepackage{booktabs,tabularx,array}
\usepackage[dvipsnames,table]{xcolor}
\usepackage{fancyhdr}
\usepackage{enumitem}
\usepackage{hyperref}

\setmainfont{Avenir Next}
\setsansfont{Avenir Next}
\setmonofont{Menlo}
\setCJKmainfont{YuMincho}
\setCJKsansfont{YuGothic}
\setCJKmonofont{YuGothic}

\definecolor{Ink}{HTML}{17253B}
\definecolor{SnowBlue}{HTML}{2F6BFF}
\definecolor{IceBlue}{HTML}{EAF2FF}
\definecolor{Glacier}{HTML}{B9D8FF}
\definecolor{SignalGold}{HTML}{F4B942}
\definecolor{RiskRed}{HTML}{C74242}
\definecolor{SafeGreen}{HTML}{16856B}
\definecolor{Muted}{HTML}{607089}
\definecolor{RuleGray}{HTML}{D8E0EB}

\hypersetup{
  colorlinks=true,
  linkcolor=SnowBlue,
  urlcolor=SnowBlue,
  pdftitle={走りやすさ指数 数理仕様},
  pdfauthor={yukisaki-center}
}

\pagestyle{fancy}
\fancyhf{}
\fancyhead[L]{\sffamily\small\color{Muted} YUKISAKI CENTER / MODEL NOTE 01}
\fancyhead[R]{\sffamily\small\color{Muted} RULE VERSION: mvp-2026-01-23-v1}
\fancyfoot[L]{\sffamily\footnotesize\color{Muted} Deterministic rule-based scoring}
\fancyfoot[R]{\sffamily\footnotesize\color{Muted} \thepage}
\renewcommand{\headrulewidth}{0.3pt}
\renewcommand{\headrule}{\hbox to\headwidth{\color{RuleGray}\leaders\hrule height \headrulewidth\hfill}}

\setlength{\parindent}{0pt}
\setlength{\parskip}{0.55em}
\setlist[itemize]{leftmargin=1.7em,itemsep=0.25em,topsep=0.25em}
\renewcommand{\arraystretch}{1.35}

\newcommand{\clip}[1]{\Pi_{[0,100]}\!\left(#1\right)}
\newcommand{\ind}[1]{\mathbf{1}\!\left[#1\right]}
\newcommand{\missing}{\varnothing}
\newcommand{\score}{\mathcal{D}}
\newcommand{\factor}[1]{\Delta_{\mathrm{#1}}}
\newcommand{\code}[1]{\texttt{#1}}

\newcommand{\callout}[2]{%
  \begin{center}
  \fcolorbox{Glacier}{IceBlue}{\begin{minipage}{0.92\linewidth}
  \vspace{0.35em}\textsf{\bfseries\color{Ink}#1}\par\smallskip
  #2
  \vspace{0.25em}\end{minipage}}
  \end{center}
}

\begin{document}

\begin{titlepage}
  \thispagestyle{empty}
  \vspace*{13mm}
  {\sffamily\small\bfseries\color{SnowBlue}\MakeUppercase{Yukisaki Center / Mathematical Specification}}\\[8mm]
  {\sffamily\fontsize{28}{36}\selectfont\bfseries\color{Ink}
    走りやすさ指数\\
    数理仕様
  }\\[4mm]
  {\sffamily\Large\color{Muted}
    Snow-road Drivability Score $\score_i(t)$
  }

  \vspace{12mm}
  {\color{SnowBlue}\rule{\linewidth}{2.4pt}}
  \vspace{9mm}

  \callout{要旨}{
    道路区間 $i$ の走りやすさを、基準点100に対する加点・減点の和として
    評価する決定的ルールモデルである。降雪、勾配、除雪経過時間、道路種別、
    消雪パイプ、凍結湿潤条件を入力し、最終値を $[0,100]$ に射影する。
    機械学習およびLLMは指数の決定に使用しない。
  }

  \vfill
  \begin{tabularx}{\linewidth}{@{}>{\sffamily\bfseries\color{Muted}}p{35mm}X@{}}
    文書版 & 1.0 \\
    対象ルール & \code{mvp-2026-01-23-v1} \\
    対象日時 & 2026年1月23日 12:00 JST \\
    対象地域 & 新潟県長岡市全域 \\
    作成日 & 2026年8月3日 \\
    実装参照 & \nolinkurl{services/drivability-scoring/src/drivability_scoring/engine.py}
  \end{tabularx}
\end{titlepage}

\section*{1. 記号と基本モデル}

道路区間 $i$、評価時刻 $t$ に対する走りやすさ指数を $\score_i(t)$ とする。
実装は初期値 $\mathcal{D}_0=100$ に各要因の寄与 $\Delta$ を加え、クリップ関数で
0以上100以下へ制限する。

\begin{equation}
\boxed{
\score_i(t)
=
\clip{
\mathcal{D}_0
+ \factor{snow,i}(t)
+ \factor{grade,i}
+ \factor{plow,i}(t)
+ \factor{road,i}
+ \factor{pipe,i}(t)
+ \factor{freeze,i}(t)
}}
\label{eq:main}
\end{equation}

\begin{equation}
\Pi_{[0,100]}(x) \coloneqq \min\!\left\{100,\,\max\{0,x\}\right\},
\qquad \mathcal{D}_0=100.
\end{equation}

ここで $\ind{A}$ は命題 $A$ が真なら1、偽なら0を返す指示関数である。
すべての寄与は整数なので、$\score_i(t)$ も整数となる。大きいほど走りやすく、
100は最も走りやすい状態、0は通行回避相当を表す。ただし、指数だけで通行可否を
決定するものではない。

\begin{table}[h]
\centering
\small
\begin{tabularx}{\linewidth}{@{}>{$}l<{$} X l@{}}
\toprule
\text{記号} & \textsf{意味} & \textsf{単位・値域} \\
\midrule
n_i(t) & 過去1時間降雪量 & cm \\
g_i & 道路区間の最大勾配 & \% \\
\tau_i(t) & 最終除雪からの経過時間 & min \\
r_i & OSM由来の道路種別 & category \\
p_i(t) & 消雪パイプが存在するか & $\{0,1\}$ \\
a_i(t) & 消雪パイプの稼働状態 & \code{active} 等 \\
\theta_i(t) & 気温 & ${}^\circ\mathrm{C}$ \\
\bottomrule
\end{tabularx}
\end{table}

\section*{2. 各要因の区分関数}

\subsection*{2.1 過去1時間降雪量 $\factor{snow}$}

\begin{equation}
\factor{snow,i}(t)=
\begin{cases}
  0,   & n_i(t)\le 0,\\
  -5,  & 0<n_i(t)<2,\\
  -15, & 2\le n_i(t)<5,\\
  -25, & 5\le n_i(t).
\end{cases}
\end{equation}

\subsection*{2.2 最大勾配 $\factor{grade}$}

\begin{equation}
\factor{grade,i}=
\begin{cases}
  0,   & g_i<5,\\
  -10, & 5\le g_i<8,\\
  -20, & 8\le g_i.
\end{cases}
\end{equation}

\subsection*{2.3 最終除雪からの経過時間 $\factor{plow}$}

除雪履歴が存在するとき、未来時刻の誤差で負の経過時間が生じないように
$\tau_i(t)$ を次式で定義する。

\begin{equation}
\tau_i(t) \coloneqq
\max\!\left\{0,\,\frac{t-t_i^{\mathrm{plow}}}{60\ \mathrm{s}}\right\}.
\end{equation}

\begin{equation}
\factor{plow,i}(t)=
\begin{cases}
  +10, & t_i^{\mathrm{plow}}\ne\missing \ \land\ 0\le\tau_i(t)\le60,\\
  -8,  & t_i^{\mathrm{plow}}\ne\missing \ \land\ 60<\tau_i(t)<180,\\
  -15, & t_i^{\mathrm{plow}}\ne\missing \ \land\ 180\le\tau_i(t),\\
  -15, & t_i^{\mathrm{plow}}=\missing.
\end{cases}
\end{equation}

履歴がない場合も「最近除雪された」と推測せず、明示的に $-15$ とする。

\subsection*{2.4 道路種別 $\factor{road}$}

集合
\begin{align}
\mathcal{R}_{15} &= \{\code{service},\ \code{track}\},\\
\mathcal{R}_{10} &= \{\code{residential},\ \code{living\_street},\ \code{unclassified}\}
\end{align}
を定め、狭幅員相当の生活道路・サービス道路を次のように評価する。

\begin{equation}
\factor{road,i}=
\begin{cases}
  -15, & r_i\in\mathcal{R}_{15},\\
  -10, & r_i\in\mathcal{R}_{10},\\
  0,   & \text{otherwise}.
\end{cases}
\end{equation}

\subsection*{2.5 稼働中の消雪パイプ $\factor{pipe}$}

\begin{equation}
\factor{pipe,i}(t)
=15\,\ind{p_i(t)=1\ \land\ a_i(t)=\code{active}}.
\end{equation}

設備が存在するだけでは加点せず、稼働状態が \code{active} の場合に限り $+15$ とする。

\subsection*{2.6 凍結湿潤条件 $\factor{freeze}$}

\begin{equation}
\factor{freeze,i}(t)
=-20\,
\ind{\theta_i(t)\ne\missing}\,
\ind{-3\le\theta_i(t)\le1}\,
\ind{n_i(t)>0}.
\end{equation}

現在の実装では、気温が $-3\,{}^\circ\mathrm{C}$ 以上 $1\,{}^\circ\mathrm{C}$ 以下で、
かつ過去1時間降雪量が正の場合を凍結湿潤条件として扱う。気温欠損時にはこの要因を
適用しない。

\clearpage
\section*{3. データ信頼度}

指数 $\score_i(t)$ と信頼度 $\mathcal{C}_i(t)$ は分離する。実装が完全性を確認する
5項目の集合を
\begin{equation}
\mathcal{K}=\{n_i,\ g_i,\ t_i^{\mathrm{plow}},\ p_i,\ \theta_i\}
\end{equation}
とすると、信頼度は次式である。

\begin{equation}
\boxed{
\mathcal{C}_i(t)
=\operatorname{round}_{2}
\left(
\frac{1}{|\mathcal{K}|}
\sum_{x\in\mathcal{K}}\ind{x\ne\missing}
\right)
}
\qquad 0\le\mathcal{C}_i(t)\le1.
\end{equation}

\code{road\_type} と \code{snow\_pipe\_operation\_status} は評価には使うが、現行の
信頼度の分母・分子には含めない。信頼度が低いことだけを理由に $\score_i(t)$ を
危険側へ変更しない。また、入力が仮データを含む場合は指数とは別に
\code{is\_simulated=true} を保持する。

\section*{4. 計算例}

\subsection*{例A：降雪中だが直近除雪・消雪設備あり}

入力を
\[
n=3,\quad g=5,\quad \tau=30,\quad p=1,\quad a=\code{active},
\quad \theta=-1
\]
とする。道路種別による減点はないものとする。このとき
\begin{align}
\score
&=\clip{100-15-10+10+0+15-20}\\
&=\boxed{80}.
\end{align}

\subsection*{例B：急勾配かつ除雪履歴なし}

入力を $n=0,\ g=9,\ t^{\mathrm{plow}}=\missing,\ p=0,\ \theta=-5$ とする。
\begin{align}
\score
&=\clip{100+0-20-15+0+0+0}\\
&=\boxed{65}.
\end{align}

\callout{クリップの意味}{
  複数のリスクが重なって合計が0未満になっても指数は0、加点後に100を超えても
  指数は100となる。したがって、出力契約 $\score\in[0,100]$ は常に維持される。
}

\clearpage
\section*{5. 実装との対応}

\begin{table}[h]
\centering
\footnotesize
\begin{tabularx}{\linewidth}{@{}l r >{\raggedright\arraybackslash}X@{}}
\toprule
\textsf{寄与} & \textsf{点数} & \textsf{出力される要因識別子} \\
\midrule
降雪 & $-5,-15,-25$ & \nolinkurl{light_hourly_snowfall},\newline \nolinkurl{moderate_hourly_snowfall},\newline \nolinkurl{heavy_hourly_snowfall} \\
勾配 & $-10,-20$ & \nolinkurl{moderate_slope}, \nolinkurl{steep_slope} \\
除雪 & $+10,-8,-15$ & \nolinkurl{plowed_within_60_minutes},\newline \nolinkurl{plowed_60_to_180_minutes_ago},\newline \nolinkurl{plowed_over_180_minutes_ago}, \nolinkurl{no_plow_history} \\
道路種別 & $-10,-15$ & \nolinkurl{narrow_road} \\
消雪パイプ & $+15$ & \nolinkurl{active_snow_pipe} \\
凍結湿潤 & $-20$ & \nolinkurl{freezing_wet_condition} \\
\bottomrule
\end{tabularx}
\end{table}

計算結果には指数・信頼度だけでなく、入力値、適用要因、評価時刻、
\nolinkurl{rule_version}、\nolinkurl{is_simulated} を保存する。S3の
\nolinkurl{curated/drivability-scores/} を正本とし、その後PostgreSQLへ投影する。

\section*{6. 適用範囲と今後の拡張}

本式は「要件上の将来像」ではなく、\textbf{現在のコードが実際に実行する式}を表す。
要件書にある次の候補は、現ルール版ではまだ $\score_i(t)$ に含まれない。

\begin{itemize}
  \item 除雪後累積降雪量、3時間降雪量、今後1時間の予測降雪量
  \item 積雪深、再凍結、橋梁、風速・吹きだまり
  \item 車線数、制限速度、標高、トンネル、通行規制
  \item 観測地点距離やデータ鮮度を使った、より詳細な信頼度補正
\end{itemize}

これらを導入する場合は、新しい寄与 $\Delta_k$ として式へ追加し、閾値・点数・
相関する要因の上限減点を定義したうえで \code{rule\_version} を更新する。
通行止めは指数の加減点ではなく、経路探索グラフから除外する制約として扱う。

\callout{安全境界}{
  この指数はルールベースの比較指標であり、道路の通行可能性や絶対的な安全を保証しない。
  LLMは確定済みの入力・要因・指数を自然言語で説明するだけで、指数、危険度、
  通行可否を決定しない。
}

\vfill
{\color{RuleGray}\rule{\linewidth}{0.5pt}}
{\sffamily\footnotesize\color{Muted}
Source of truth: \nolinkurl{services/drivability-scoring/src/drivability_scoring/engine.py}\\
Rule version: \code{mvp-2026-01-23-v1}
}

\end{document}
